\section*{序\qquad 言}
\addcontentsline{toc}{section}{序言}
\label{sec:b4_preface}

第三冊第一章一直在指示三、四這兩冊書的走向．

我們說過數學是在解決實際問題的刺激與要求下才發展數學的．但這句話須進一層的了解，事實上許多數學發展「主要的產物」並不局限於只能處理一個特殊而枝節的問題，毋寧說它發展成功的是一種方法，一種能一併解決整堆同一類型問題的系統性的方法．初中課程裏引入代數方程式的求解來替代算術，第三冊笛卡兒座標幾何及其後發展出來的向量幾何來替代平面幾何都是沿依著這一方向的主要發展．

在這種意味下，數學發展與實際問題的配合是整套的配合，而不是平常意義下的「應用」所能概括的．

而數學教育的內容也應以傳授這種主產物為主，這個精神撐起了實驗教材的兩大支柱：
\begin{enumerate}[(i)]
  \item 尊重數學發展的原意，使不與實際問題脫節．
  \item 專注系統性的方法．
\end{enumerate}

就 (i) 說來，到第四冊我們的數學工具已經逐漸完備．在各個工具的製作中，我們插入一些實際問題的解法．這些實際問題包含自然與人文的應用．由於牽涉一些其他學科的專門知識，我們能收集在本冊的還很零星不全，但相信多少仍能見其一斑．

就 (ii) 來說，我們所謂的數學工具，指的是系統性的方法，而不是在「嚴密與系統化」的美名下作數學名詞的玩弄．這系統性的方法繞著兩個主題：
\begin{enumerate}[(一)]
  \item \textbf{運算工具本身的發展}：這包含了方程式的求解（如第一冊），數的觀念的澄清（如第二冊），以及各種基本函數的討論（如第一、二冊：多項式函數，對數函數，指數函數，第三冊：三角函數）．
  \item \textbf{把問題經由幾何系統性地轉化成運算對象}：這一精神在第三冊的向量與座標幾何中多少已經交代了．
\end{enumerate}

其實我們一直是在揉合這兩個主題．我們提過（第三冊第一章）自然或人文的現象常以函數的姿態出現，但這些函數可以千變萬化，我們必須尋找一些比較簡單的函數作為基本的模型，來逼近地描寫這些一般時候可能是奇形怪狀的函數，例如我們已多少把握了一次函數，二次函數，對數與指數函數及三角函數了．

第三冊我們尤其強調一次函數及其所對應的一次方程式，並把它們所對應的直線，平面等的幾何性質的討論統稱為線性幾何，那裏所標明的向量幾何與座標幾何兩章嚴格地說來都只屬於線性幾何的範圍的．

這一冊我們要開始精細地討論二次函數，看看二次方程式所相應的幾何對象是什麼？不過在此之前我們還要考慮“動態的”線性幾何，卽考慮所謂線性映射及其相關的性質．因為映射的觀點（亦卽不變性的討論）在數學與物理的發展上都扮演著極端重要的角色．

複數及高斯平面也要同時引入，用複數來描寫平面上的幾何．在代數的發展上這會是一個重要的關鍵．

至於傳統的三角學，就要在「三角學的基本定理」的出現之後結束它雄踞教壇數十年的歷程．這是複數的功用，也是我們系統性方法的一面．

\vspace{2em}
\begin{flushright}
\textbf{高中數學實驗教材編輯小組}
\end{flushright}
